\section*{Hoofdstuk \ref{ch:b3:endocrinology} --- Endocrinologie en homeostase}

\begin{solution}{exo:b3:endocrinology:1}
\omterm{def:b3:endocrinology:glucose}{Insuline}: peptide, receptor aan het oppervlak (een tyrosinekinase),
werkt in minuten, halveringstijd ongeveer \qty{5}{min}. Cortisol:
steroïde, \omterm{def:b3:endocrinology:hormone}{kernreceptor}, uren, halveringstijd \qty{80}{min} (grotendeels
gebonden aan een transporteiwit). Adrenaline: amine, receptoren aan het
oppervlak gekoppeld aan een G-eiwit, seconden, halveringstijd ongeveer
\qty{2}{min}. Thyroxine: gejodeerd aminozuur, \omterm{def:b3:endocrinology:hormone}{kernreceptor}, dagen,
halveringstijd \qty{7}{dagen} (gebonden aan transporteiwitten).
Vasopressine: peptide, receptor aan het oppervlak gekoppeld aan een
G-eiwit (cAMP in de \omterm{def:b3:renal-osmoregulation:nephron}{verzamelbuis}), minuten, halveringstijd ongeveer
\qty{15}{min}.
\end{solution}

\begin{solution}{exo:b3:endocrinology:2}
\omterm{def:b3:endocrinology:axes}{Hypothalamus} (TRH) $\to$ \omterm{def:b3:endocrinology:axes}{hypofyse} (TSH) $\to$ schildklier (T$_{4}$),
waarbij T$_{4}$ beide bovenste lagen remt. Vernietigde schildklier:
T$_{4}$ laag, TSH hoog. \omterm{def:b3:cancer-biology:hallmarks}{Tumor} die TSH uitscheidt: TSH hoog en T$_{4}$
hoog (de terugkoppeling onderdrukt de \omterm{def:b3:cancer-biology:hallmarks}{tumor} niet). Te veel
thyroxinetabletten: T$_{4}$ hoog, TSH onderdrukt. Jodiumgebrek:
T$_{4}$ laag of laag-normaal, TSH hoog, en de klier groeit eronder uit
tot een struma.
\end{solution}

\begin{solution}{exo:b3:endocrinology:3}
Het brein sterft binnen minuten aan een lage glucose en wordt door een
hoge pas na jaren geschaad; en de evolutie kwam de honger veel vaker
tegen dan het feestmaal. Dus is de verdediging tegen een daling
overtollig aanwezig (\omterm{def:b3:endocrinology:glucose}{glucagon}, adrenaline, cortisol, groeihormoon) en
de verdediging tegen een stijging enkelvoudig. Voorspelling: te veel
\omterm{def:b3:endocrinology:glucose}{insuline} is acuut dodelijk (hypoglykemisch coma), te weinig is een
chronische ziekte --- en dat is wat de kliniek ziet.
\end{solution}

\begin{solution}{exo:b3:endocrinology:4}
Een \omterm{def:b3:endocrinology:clock}{zeitgeber} is een aanwijzing uit de omgeving die een klok
gelijkzet: licht, het tijdstip van het eten, temperatuur, beweging, het
sociale rooster. De hoofdklok gebruikt licht, via de
melanopsine-ganglioncellen van het \omterm{def:b3:sensory-systems:retina}{netvlies}. Na acht tijdzones
verschuift de klok maar ongeveer een uur per dag, zodat slaap,
cortisol, temperatuur en eetlust een week lang op de oude tijd lopen
terwijl de perifere klokken, door de maaltijden bijgesteld, in hun
eigen tempo meedrijven: jetlag, erger oostwaarts, omdat de klok
vervroegen moeilijker is dan hem vertragen.
\end{solution}

\begin{solution}{exo:b3:endocrinology:5}
Half maximaal antwoord bij $250$ bezette receptoren: $H = K_{d}\times
250/(\num{20000} - 250) = \qty{1}{nM}\times 0.0127 = \qty{12.7}{pM}$,
tachtig maal onder $K_{d}$. Met \num{2000} receptoren: $250/1750
\times \qty{1}{nM} = \qty{143}{pM}$, elf maal minder gevoelig. Een
aanhoudend hoge \omterm{def:b3:endocrinology:glucose}{insuline} reguleert haar receptoren neer, zodat een
gegeven effect meer \omterm{def:b3:endocrinology:glucose}{insuline} vergt: één onderdeel van de
ongevoeligheid voor \omterm{def:b3:endocrinology:glucose}{insuline}, die zichzelf voedt.
\end{solution}

\begin{solution}{exo:b3:endocrinology:6}
$T_{4} = 13$: $1.5\,\mathrm{e}^{0.92} = \qty{3.8}{mU/L}$; $11$:
$1.5\,\mathrm{e}^{1.84} = \qty{9.4}{mU/L}$; $9$: $1.5\,\mathrm{e}^{2.76}
= \qty{24}{mU/L}$; $25$: $1.5\,\mathrm{e}^{-4.6} = \qty{0.015}{mU/L}$.
Bij \qty{11}{pmol/L} ligt het T$_{4}$ nog binnen \qtyrange{10}{20}{}
terwijl de TSH meer dan het dubbele van haar bovengrens is
(subklinische hypothyreoïdie); bij $13$ valt de TSH nog net binnen het
gebied; bij $9$ zijn beide afwijkend; bij $25$ zijn beide de andere
kant op afwijkend.
\end{solution}

\begin{solution}{exo:b3:endocrinology:7}
$L = s\beta/(a\gamma) = 0.2\times 0.05/(0.02\times 0.1) = 5$.
Stationaire overmaat $g^{*} = (u/a)/(1 + L) = (2/0.02)/6 = \qty{16.7}{mg/dL}$;
zonder terugkoppeling $u/a = \qty{100}{mg/dL}$. Overmaat aan \omterm{def:b3:endocrinology:glucose}{insuline}
$i^{*} = \beta g^{*}/
\gamma = 0.05\times 16.7/0.1 = \qty{8.3}{mU/L}$.
\end{solution}

\begin{solution}{exo:b3:endocrinology:8}
Discriminant $(a - \gamma)^{2} - 4s\beta = 0.0064 - 0.04 < 0$:
oscillerend. $\omega = \sqrt{a\gamma + s\beta - (a+\gamma)^{2}/4} =
\sqrt{0.002 + 0.01 - 0.0036} = \qty{0.092}{min^{-1}}$, periode $2\pi/
\omega = \qty{68}{min}$; de amplitude daalt met een factor $\mathrm{e}$
in $2/(a +
\gamma) = \qty{16.7}{min}$. Halvering van $s$: $s\beta = 0.005$,
discriminant $0.0064 - 0.02 < 0$, nog altijd oscillerend,
$\omega = \sqrt{0.007 - 0.0036}
= \qty{0.058}{min^{-1}}$, periode \qty{108}{min}; de vervaltijd blijft
gelijk; $L = 2.5$. Een zwakkere lus keert trager terug en houdt een
grotere blijvende fout.
\end{solution}

\begin{solution}{exo:b3:endocrinology:9}
Een jaar lang onderdrukte de prednisolon CRH en ACTH; de bijnierschors,
niet geprikkeld, atrofieerde. Bij het stoppen is er geen steroïde van
buitenaf meer, hebben \omterm{def:b3:endocrinology:axes}{hypothalamus} en \omterm{def:b3:endocrinology:axes}{hypofyse} weken nodig om te
hervatten, en zou de geslonken bijnier niet op ACTH kunnen antwoorden
als dat kwam. Een besmetting vraagt tienmaal de normale hoeveelheid
cortisol; is die er niet, dan verwijden de vaten en dalen druk en
glucose: een bijniercrisis. Gemeten: ACTH laag (de laesie zit
centraal), cortisol laag, en een geringe stijging van cortisol na
ingespoten ACTH (de klier is verschrompeld) --- anders dan bij de
ziekte van Addison, waar de ACTH hoog is. Vandaar het langzame
afbouwen.
\end{solution}

\begin{solution}{exo:b3:endocrinology:10}
Eén variabele: $\dot x = f(x)$ met $f' < 0$ heeft één vast punt
$x^{*}$; voor $x > x^{*}$ is $\dot x < 0$ en daalt $x$ monotoon naar
$x^{*}$ zonder het ooit te passeren (de eenduidigheid van de
oplossingen), en onder het punt symmetrisch: geen trilling. Twee
variabelen: de matrix kan complexe eigenwaarden hebben, zoals in de
stelling, wanneer de uitvoerder via een tweede variabele met haar eigen
tijdschaal werkt. Een onmiddellijke onderdrukking zou op een vaste
waarde uitkomen; de vertraging tussen transcriptie en onderdrukking in
de kern (translatie, dimerisatie, fosforylering, invoer) laat het eiwit
doorschieten voordat het gen wordt uitgezet en te ver dalen voordat het
weer wordt aangezet, en de periode is ruwweg tweemaal de som van de
vertraging en de tijd om het eiwit af te breken. Een snellere afbraak
van PER verkort de periode: een destabiliserende mutatie van het
menselijke PER2 geeft een klok van ongeveer 20 uur en een familie die
om zeven uur 's avonds in slaap valt.
\end{solution}

\begin{solution}{exo:b3:endocrinology:11}
Exponent $k(\text{Ca} - 1.2)$: bij $1.10$, $\mathrm{e}^{-2}$, PTH $= 1/
(1 + 0.135) = 0.88$; $1.15$: $0.73$; $1.20$: $0.50$; $1.25$: $1/(1 +
2.72) = 0.27$; $1.30$: $1/(1 + 7.39) = 0.12$. Helling in het \omterm{thm:b3:endocrinology:loop}{instelpunt}
$-k/4 = -5$ per mmol/L: \qty{5}{\%} van het maximum voor elke
\qty{0.01}{mmol/L}. Calcium moet binnen enkele procenten worden
gehouden en zijn uitvoerder werkt op een enorme buffer (het bot) zonder
door te schieten, dus is een steil antwoord veilig en nodig. Een
glucoselus die zo steil was, werkend via \omterm{def:b3:endocrinology:glucose}{insuline} met haar eigen
klaringstijd, zou een grote \omterm{thm:b3:endocrinology:loop}{lusversterking} hebben en na elke maaltijd
diep de hypoglykemie in oscilleren.
\end{solution}

\begin{solution}{exo:b3:endocrinology:12}
Glycering is onomkeerbaar en traag, dus heeft een cel van leeftijd
$\tau$ een fractie $kG\tau$ geglyceerd; gemiddeld over cellen die
gelijkmatig \numrange{0}{120} dagen oud zijn, is HbA1c
$= kG\times\qty{60}{dagen}$, evenredig met de gemiddelde glucose.
IJking: $5.5\times 60k = 0.05$ geeft $k =
\qty{1.5e-4}{}$ per (mmol/L)(dag); bij \qty{10}{mmol/L}: \qty{9.1}{\%}.
Met cellen die 60 dagen leven is de gemiddelde leeftijd 30 dagen en is
HbA1c bij dezelfde glucose gehalveerd: een patiënt bij
\qty{10}{mmol/L} zou \qty{4.5}{\%} aflezen, normaal.
\end{solution}

\begin{solution}{pb:b3:endocrinology:1}
\textbf{1.} \qty{90}{mg/dL} $= \qty{0.9}{g/L} = 0.9/180 = \qty{5.0}{mmol/L}$;
$0.9\times 15 = \qty{13.5}{g}$ in het verdelingsvolume.
\textbf{2.} $75/15 = \qty{5}{g/L}$: een stijging van \qty{500}{mg/dL},
\qty{28}{mmol/L}. De werkelijke piek ligt veel lager omdat de opname
over een uur is uitgesmeerd terwijl de afvoer de glucose weghaalt
zodra zij aankomt, en de lever bij de eerste passage een groot deel
opneemt.
\textbf{3.} $\qty{75}{g}/\qty{60}{min} = \qty{1.25}{g/min}$; over
\qty{15}{L}: \qty{8.3}{mg/dL/min}, viermaal $u_{0}$.
\textbf{4.} Stationaire overmaat $u/a = 8.3/0.02 = \qty{417}{mg/dL}$,
tijdconstante $1/a = \qty{50}{min}$; na een uur $417(1 -
\mathrm{e}^{-1.2}) = \qty{291}{mg/dL}$ boven nuchter: ongeveer
\qty{380}{mg/dL}, \qty{21}{mmol/L}.
\textbf{5.} $L = 0.2\times 0.05/(0.02\times 0.1) = 5$; stationaire
overmaat $417/6 = \qty{69}{mg/dL}$, zes maal kleiner dan zonder
\omterm{def:b3:endocrinology:glucose}{insuline}.
\textbf{6.} Het brein verbrandt \qty{120}{g} glucose per dag, slaat
niets op en kan weinig anders gebruiken: een lage glucose geeft binnen
minuten verwardheid en coma. Een hoge glucose glyceert eiwitten en
beschadigt in de loop van jaren de kleine vaten van \omterm{def:b3:sensory-systems:retina}{netvlies}, nier en
zenuwen en de grote slagaders.
\textbf{7.} $\lambda^{2} + (a + \gamma)\lambda + (a\gamma + s\beta) = 0$:
spoor $-0.12$, determinant $0.002 + 0.01 = 0.012$, discriminant
$0.0144 - 0.048 = -0.0336$.
\textbf{8.} $\omega = \sqrt{0.0336}/2 = \qty{0.092}{min^{-1}}$; periode
\qty{68}{min}; vervaltijd $2/0.12 = \qty{16.7}{min}$.
\textbf{9.} $\dot g(0) = 100(-0.06 + c\,\omega) = -2$, dus $c\,\omega =
0.04$, $c = 0.44$.
\textbf{10.} $g = 0$ wanneer $\tan\omega t = -1/c = -2.29$, $\omega t =
\pi - 1.16 = 1.98$, $t = \qty{21.6}{min}$. Bij het eerste minimum, rond
\qty{32}{min}: $100\,\mathrm{e}^{-1.94}(\cos 2.97 + 0.44\sin 2.97) =
14.4\times(-0.91) \approx -\qty{13}{mg/dL}$, glucose \qty{77}{mg/dL}.
\textbf{11.} Bij $t = 0$ is $\dot i = 100\beta > 0$: de \omterm{def:b3:endocrinology:glucose}{insuline} stijgt
zolang $g > 0$, en piekt wanneer $\beta g = \gamma i$, wat vergt dat
$g$ nog positief is --- dus komt de piek voordat de glucose nuchter
bereikt, en nadat de glucose is beginnen te dalen (zij daalt vanaf
$t = 0$). In het model ligt de piek rond \qty{12}{min}.
\textbf{12.} Met invoer die over een uur is uitgesmeerd valt de piek
wanneer opname en afvoer in evenwicht zijn, lager en later (30--60
min); de terugkeer wacht tot de opname is afgelopen, vandaar twee uur;
het doorschot van de \omterm{def:b3:endocrinology:glucose}{insuline} is bij een geleidelijke invoer kleiner,
dus is de dip mild.
\textbf{13.} $L = 2.5$. Gezond $g^{*} = (2/0.02)/6 = \qty{16.7}{mg/dL}$;
ongevoelig $100/3.5 = \qty{28.6}{mg/dL}$: een stijging van ongeveer
\qty{12}{mg/dL} (\qty{0.7}{mmol/L}), nuchtere glucose rond
\qty{102}{mg/dL}, \qty{5.7}{mmol/L} --- de drempel van een gestoorde
nuchtere glucose.
\textbf{14.} $i^{*} = \beta g^{*}/\gamma$: gezond $8.3$, ongevoelig
\qty{14.3}{mU/L}, verhouding $1.7$: gecompenseerde ongevoeligheid voor
\omterm{def:b3:endocrinology:glucose}{insuline}, hyperinsulinemie bij een bijna normale glucose.
\textbf{15.} $L = 0.1\times 0.01/0.002 = 0.5$; $g^{*} = 100/1.5 =
\qty{66.7}{mg/dL}$, \qty{50}{mg/dL} boven de gezonde toestand:
\qty{2.8}{mmol/L}, nuchtere glucose \qty{7.8}{mmol/L}
(\qty{140}{mg/dL}), boven de diabetesdrempel.
\textbf{16.} Discriminant $(a - \gamma)^{2} - 4s\beta = 0.0064 - 0.004 =
0.0024 > 0$: monotone terugkeer. $\lambda = -0.06 \pm 0.0245$: $-0.0355$
en $-0.0845$ per minuut; traagste tijdconstante \qty{28}{min}, tegen
\qty{16.7}{min} voor het gezonde verval: de falende lus keert trager
terug en schiet nooit door.
\textbf{17.} Hoge glucose bij hoge \omterm{def:b3:endocrinology:glucose}{insuline}: ongevoeligheid met
gedeeltelijke opvang --- type~2; de waarde na twee uur van
\qty{13}{mmol/L} overschrijdt \qty{11.1}{}, dus diabetes en niet
slechts een gestoorde tolerantie.
\textbf{18.} Geen C-peptide betekent geen eigen \omterm{def:b3:endocrinology:glucose}{insuline}: type~1.
Zonder \omterm{def:b3:endocrinology:glucose}{insuline} kunnen spier en vet geen glucose opnemen, wordt vet
afgebroken tot vetzuren en ketonen, wordt spiereiwit afgebroken voor de
gluconeogenese, en gaat glucose met haar calorieën in de urine
verloren: de patiënt verhongert te midden van overvloed.
\textbf{19.} $5\times 9/5.5 = \qty{8.2}{\%}$.
\textbf{20.} $1.5\,\mathrm{e}^{0.46\times 3} = 1.5\times 3.97 =
\qty{6.0}{mU/L}$. T$_{4}$ binnen het gebied, TSH erboven: subklinische
hypothyreoïdie, met een falende klier en een compenserende \omterm{def:b3:endocrinology:axes}{hypofyse}.
\textbf{21.} $1.5\,\mathrm{e}^{0.46\times 8} = 1.5\times 39.6 =
\qty{59}{mU/L}$. Een laag T$_{4}$ bij een TSH van slechts $0.3$
betekent dat de \omterm{def:b3:endocrinology:axes}{hypofyse} niet antwoordt: de laesie zit centraal
(\omterm{def:b3:endocrinology:axes}{hypofyse} of \omterm{def:b3:endocrinology:axes}{hypothalamus}), niet in de schildklier.
\textbf{22.} Eén halveringstijd: de spiegel daalt tot \qty{50}{\%}, en
de klachten sluipen er in de loop van weken in. Cortisol, met een
halveringstijd van \qty{80}{min}, is binnen uren na een gemiste dosis
verdwenen, en een belasting op die dag treft geen verdediging aan: de
vergeten week thyroxine is onaangenaam, de vergeten dag cortisol kan
doden.
\textbf{23.} T$_{4}$ hoog (het \omterm{def:b3:adaptive-immunity:antibody}{antilichaam} drijft de klier aan ongeacht
de terugkoppeling); TSH onderdrukt door het hoge T$_{4}$; het
log-lineaire verband drijft haar tot onaantoonbaar --- bij $T_{4} = 30$
is $1.5\,\mathrm{e}^{-6.9}
= \qty{0.0015}{mU/L}$ --- zodat een onaantoonbare TSH het eerste teken
is.
\textbf{24.} De spiegel van een \omterm{def:b3:endocrinology:hormone}{hormoon} meldt het evenwicht tussen zijn
bevel en zijn klier: een hoge TSH bij een laag T$_{4}$ is een falende
klier en een hoge TSH bij een hoog T$_{4}$ een op hol geslagen
\omterm{def:b3:endocrinology:axes}{hypofyse}; een hoge \omterm{def:b3:endocrinology:glucose}{insuline} bij een hoge glucose is een ongevoelig
lichaam en een lage \omterm{def:b3:endocrinology:glucose}{insuline} bij een hoge glucose een vernietigd
eilandje. Alleen gelezen zegt een T$_{4}$ van \qty{7}{pmol/L} of een
\omterm{def:b3:endocrinology:glucose}{insuline} van \qty{30}{mU/L} niet waar de fout zit; gelezen met zijn
bevelvoerder wel.
\textbf{25.} \omterm{thm:b3:endocrinology:loop}{Lusversterking} $L = 5$; periode \qty{68}{min} en vervaltijd
\qty{17}{min}; nuchtere glucose van de falende lus \qty{7.8}{mmol/L}.
\end{solution}
